\documentclass[12pt,a4paper,UTF8]{ctexart}

%==========================================================================
% 宏包
%==========================================================================
\usepackage{geometry}
\usepackage{amsmath,amssymb}
\usepackage{graphicx}
\usepackage{booktabs}
\usepackage{hyperref}
\usepackage{float}
\usepackage{enumitem}
\usepackage{multirow,array}
\usepackage{cite}
\usepackage{caption}

%==========================================================================
% 页面布局
%==========================================================================
\geometry{left=3cm,right=2.5cm,top=2.5cm,bottom=2.5cm}
\hypersetup{colorlinks=true,linkcolor=blue,citecolor=blue,urlcolor=blue}

%==========================================================================
% 标题
%==========================================================================
\title{\textbf{基于耦合多智能体强化学习的\\低空交叉路口飞行器协同管控模型}}
\author{彭义森\quad 指导老师：韩云祥\\[4pt]\small 四川大学}
\date{2026年7月}

\begin{document}
\maketitle

%==========================================================================
% 摘要
%==========================================================================
\begin{abstract}
城市低空交叉路口场景中电动垂直起降飞行器（eVTOL）的高密度安全协同管控是先进空中交通（AAM）面临的核心挑战之一。本文提出一种基于动态图注意力网络（DyGAT）的耦合多智能体强化学习（MARL）模型，用于解决该场景下的冲突避免与流量协同优化问题。模型包含三项核心设计：（1）DyGAT融合空间注意力、时序注意力与基于物理量（距离、速度差、航向差）的连续动态边权重，刻画飞行器间时空耦合关系；（2）分层奖励函数与动态安全距离，兼顾安全与效率；（3）课程学习与基于冲突事件的episode级损失加权（$\lambda=2.0$），在保持PPO on-policy框架的前提下提升稀疏冲突样本的训练效率。

在21架eVTOL的二维交叉路口仿真场景中，本文方法与八种基线方法（涵盖规则调度、单智能体RL、独立MARL、CTDE-MARL及工程安全方法）进行了系统对比。相比CTDE-MARL中性能最强的MAPPO方法，本文方法的总冲突次数（TC）从$20.3\pm2.4$降至$18.2\pm2.1$（降幅10.3\%，$p=0.180$，Cohen's $d=0.93$），无冲突完成率（CFCR）从$57.8\%$提升至$68.4\%$（$p<0.001$，Cohen's $d=5.27$），最小间隔违反率（MSVR）从$3.0\%$降至$2.4\%$（$p=0.029$）。在工程安全方法中，ORCA的TC为22.8、CFCR约50\%，验证了学习方法在高密度场景下的协同优势。

本文在讨论部分坦诚分析了方法的局限性，包括二维场景简化、$n=5$小样本量导致的TC统计检验力不足（约38\%）、均匀风扰模型的简化处理，以及eVTOL加速度参数缺乏制造商公开规格的问题。

\noindent\textbf{关键词：}低空交通管控；多智能体强化学习；动态图注意力网络；冲突避免；集中训练分布式执行；无冲突完成率
\end{abstract}

%==========================================================================
% 1. 引言
%==========================================================================
\section{引言}

\subsection{研究背景与动机}

城市地面交通拥堵的加剧推动了以eVTOL为核心的城市空中交通（UAM）发展。FAA预测到2040年将有数百万架无人机与eVTOL在美国低空空域运行\cite{ref1}。与传统民航航路的结构化交通不同，城市低空交通具有高密度、强动态特征——飞行器在空间中快速移动，邻域拓扑以秒级频率变化。在此背景下，交叉路口与汇合点的协同管控与冲突避免是低空无人机交通管理（UTM）面临的核心挑战之一。

传统空中交通流量管理（ATFM）方法包括整数规划、蒙特卡洛仿真与基于欧拉模型的线性规划等\cite{ref2}。这些方法在可预测场景中效果良好，但计算复杂度随飞行器数量指数增长。模型预测控制（MPC）\cite{ref3}、最优互惠碰撞避免（ORCA）\cite{ref4}与控制屏障函数（CBF）\cite{ref5}等工程方法在形式化安全保证方面具有优势，但在高密度交互场景中趋于保守。

近年来，MARL被成功引入交通管控领域。Agogino与Tumer\cite{ref6,ref7}建立了分布式航路点智能体调控框架；Brittain与Wei\cite{ref8}在BlueSky仿真器中验证了深度MARL的冲突规避能力；Ghosh等\cite{ref9}提出深度集成MARL方法提升决策鲁棒性；Li等\cite{ref10}提出MS-GRL采用多尺度图增强RL解决密集无人机网络冲突。然而，现有工作在三个维度上仍存在不足：

\textbf{（1）时空耦合建模不足。}多数方法使用全连接网络或静态图结构编码智能体间关联。当飞行器快速移动导致邻域拓扑剧烈变化时，静态图结构无法区分"距离相同但相对运动状态不同"的飞行器对的差异化风险——两架相向飞行的飞行器与两架同向飞行的飞行器，即便距离相等，冲突风险相差数个数量级。

\textbf{（2）稀疏冲突样本利用不足。}低空交通中冲突属低概率高风险事件，训练中冲突样本占比通常低于$0.1\%$。标准随机采样的mini-batch训练中，冲突样本极易被海量正常样本淹没\cite{ref11,ref12}。

\textbf{（3）环境不确定性建模不足。}多数研究采用简化模型，忽略风扰、传感器噪声与物理运动约束\cite{ref13}。此外，现有研究多采用固定安全距离，未考虑不同飞行速度下的差异化制动需求。

针对上述不足，本文设计了一种耦合MARL管控模型：通过DyGAT的物理量驱动动态边权重刻画时空耦合关系；通过课程学习与episode级损失加权高效利用稀疏冲突样本；通过动态安全距离与物理约束提升环境建模的工程合理性。

\subsection{主要贡献}

\begin{enumerate}[label=(\arabic*)]
    \item 构建了面向低空交叉路口冲突规避的仿真环境，引入与速度正相关的动态安全距离模型与eVTOL物理运动约束；
    \item 提出DyGAT网络，通过空间注意力、时序注意力与基于距离、速度差、航向差的连续动态边权重，刻画飞行器间细粒度时空耦合关系；
    \item 设计课程学习与episode级冲突损失加权训练框架（$\lambda=2.0$），在保持PPO on-policy理论框架的前提下提升稀疏冲突样本的训练效率；
    \item 构建双层评价指标体系——区分表层指标（CR、TC、AD）与深层指标（CFCR、MCAE、MSVR），揭示通行完成率（99.1\%）与无冲突完成率（68.4\%）之间的显著差距；
    \item 与八种基线方法（涵盖规则调度RR/FCFS、单智能体DQN、独立A2C、CTDE-MARL方法MADDPG/COMA/QMIX/MAPPO，以及工程安全方法ORCA）进行了系统对比。
\end{enumerate}

%==========================================================================
% 2. 相关工作
%==========================================================================
\section{相关工作}

\subsection{传统ATFM与工程安全方法}

传统ATFM方法可分为集中式优化与分布式启发式两类。集中式优化包括整数规划和基于欧拉模型的线性规划\cite{ref2}，能够提供最优性保证但面临可扩展性挑战。启发式方法中，Agogino与Tumer\cite{ref14}设计了基于规则的反馈控制机制，计算效率高但泛化能力有限。MPC\cite{ref3}通过滚动时域优化实现预测性控制；ORCA\cite{ref4}通过速度空间的几何约束实现分布式碰撞避免，计算效率高且具有形式化安全保证；CBF\cite{ref5}通过Lyapunov类约束将安全需求编码为控制输入限制。

\subsection{强化学习在交通管控中的应用}

单智能体RL方面，DQN已被应用于单飞行器冲突规避与交叉口信号控制\cite{ref11}。国内研究者在交通信号控制的RL应用方面也开展了大量工作\cite{ref15,ref16,ref17,ref18,ref19,ref20,ref21,ref22}。

多智能体RL方面，Ghavamzadeh等\cite{ref23}提出层级式MARL框架；Tumer与Agogino\cite{ref7,ref14,ref24}系统验证了分布式架构在ATM中的有效性；Ghosh等\cite{ref9}提出深度集成MARL方法；翟子洋等\cite{ref12}系统梳理了从独立Q-learning到QMIX、MAPPO等CTDE方法的演进路线。

Li等\cite{ref10}提出的MS-GRL与本文目标最为接近——同样面向密集低空网络的冲突解决，同样采用图注意力机制。两者在方法学上的主要差异为：MS-GRL采用Double DQN处理离散速度档位，本文采用PPO处理连续加速度控制；MS-GRL的图结构基于固定距离阈值，本文的DyGAT通过物理量驱动的连续动态边权重区分冲突风险等级；MS-GRL面向100架UAV通用场景，本文聚焦21架eVTOL交叉路口场景。

\subsection{地面交通信号控制的RL研究}

地面交通信号控制领域的RL研究为本文提供了方法论参考。Aslani等\cite{ref11}验证了actor-critic架构在连续动作空间中的优势；赖建辉\cite{ref17}探索了优先经验回放在交通控制中的应用；翟子洋等\cite{ref12}指出CTDE架构在多路口协同中的潜力。尽管低空交通与地面交通在物理约束上存在本质差异，上述工作在训练机制和评估方法论上的经验具有可迁移性。

%==========================================================================
% 3. 问题建模与方法
%==========================================================================
\section{问题建模与管控模型设计}

本文聚焦城市低空十字交叉路口场景：$N$架eVTOL从均匀分布于圆周的不同方向飞向交叉区域中心，需在保持安全间隔的前提下以最小延迟通过交叉区域。

\subsection{分布式马尔可夫决策过程形式化}

将$N$架eVTOL的协同管控形式化定义为Dec-MDP六元组$\langle \mathcal{N}, \mathcal{S}, \mathcal{A}, \mathcal{P}, \mathcal{R}, \gamma \rangle$：

\begin{itemize}
    \item \textbf{智能体集合} $\mathcal{N} = \{1, 2, \ldots, N\}$：每架eVTOL为一个独立决策智能体。训练阶段$N$由课程学习动态调整（$5\to21$），测试阶段$N=21$。
    \item \textbf{全局状态} $\mathcal{S}$：$S = (s_1, \ldots, s_N)$，其中$s_i = [x_i, y_i, v_i, v_i^{\text{target}}, d_i^{\text{center}}]$为飞行器$i$的五维形式化状态（二维位置、当前速度、目标速度、距中心距离）。需要说明：该五维向量是Dec-MDP的形式化定义；智能体网络的实际输入是经邻域扩展与历史序列编码后的局部观测$o_i$（见第3.2.1节），全局状态$S$仅在CTDE训练阶段供GAT-Mixer计算全局价值函数使用。
    \item \textbf{联合动作} $\mathcal{A}$：$A = (a_1, \ldots, a_N)$，$a_i \in \mathbb{R}$为连续加速度（m/s$^2$），经运动学约束与风扰影响后产生实际速度变化。
    \item \textbf{状态转移} $\mathcal{P}$：受飞行器运动学模型（第3.2.2节）、风扰与传感器噪声（第3.2.3节）的共同影响。
    \item \textbf{联合奖励} $\mathcal{R}$：全局奖励$R = \frac{1}{N}\sum_{i=1}^{N} r_i$，$r_i$由四个分量构成（第3.3节）。
    \item \textbf{折扣因子} $\gamma = 0.99$。
\end{itemize}

\subsection{环境模型}

仿真环境以二维平面模拟交叉路口空域（$1000\text{m} \times 1000\text{m}$）。$N$架飞行器从圆周均匀分布的初始位置出发，航向指向中心区域。

\subsubsection{观测空间}

为避免概念混淆，本文将状态/观测信息分为三个层次：

\textbf{第一层——形式化状态} $s_i \in \mathbb{R}^5$：$[x_i, y_i, v_i, v_i^{\text{target}}, d_i^{\text{center}}]$。这是Dec-MDP的形式化定义，仅描述单个飞行器的基础信息。

\textbf{第二层——局部观测} $o_i$：以$s_i$为基础进行空间扩展（邻域飞行器的相对位置、相对速度、加速度与航向角）和时间扩展（最近$T=10$步历史序列）。这是智能体网络的直接输入。

\textbf{第三层——编码特征} $h_i^{\text{coupled}} \in \mathbb{R}^{256}$：$o_i$经LSTM编码和DyGAT聚合后的256维耦合特征向量。数据流为：$s_i \to o_i \to \text{LSTM} \to \text{DyGAT} \to h_i^{\text{coupled}} \to$ 策略/Q值输出。

\subsubsection{动作空间与物理运动约束}

动作空间为连续加速度$a_i \in [-4.0, 3.0]$ m/s$^2$（最大减速度4.0 m/s$^2$，最大加速度3.0 m/s$^2$）。需要诚实指出：当前主流eVTOL制造商（Joby、Volocopter、Lilium等）均未公开最大水平加速度的具体技术规格。本文的加速度取值参考了学术文献中对eVTOL的典型建模假设\cite{ref25}——该文献在eVTOL性能比较研究中使用2.0 m/s$^2$作为标准加速度假设——并在此基础上考虑了一定的安全裕度（紧急制动时减速度高于正常加速度），将其作为仿真环境中的代表性参数设定，而非对特定机型的精确复现。

速度范围设为$[60, 140]$ km/h，参考了当前eVTOL的技术水平（如Volocopter约90--110 km/h，Joby S4巡航约240 km/h，EHang 216约100--130 km/h）。最大转弯角设为$\pi/6$（30$^\circ$）。速度更新模型为：

\begin{equation}
    v_i^{t+1} = \text{clip}\left(v_i^{t} + a_i \cdot \Delta t \cdot 3.6,\; 60,\; 140\right)
\end{equation}

其中$\Delta t = 0.5$ s为仿真步长，系数3.6实现m/s$^2$到km/h的单位转换。

\subsubsection{不确定性建模与动态安全距离}

\textbf{传感器噪声：}对位置观测添加$\mathcal{N}(0, 0.5^2)$高斯噪声（m），仅作用于观测通道。

\textbf{风扰建模：}为每架飞行器逐时间步添加随机方向$\theta_w \sim \mathcal{U}(0, 2\pi)$、固定强度$w = 0.1$ m/s的风扰速度矢量。该模型是对真实低空风场的简化——真实风场通常使用Dryden（MIL-F-8785C）\cite{ref26}或Kaimal（IEC 61400-1）谱模型刻画，具有空间相关性与时间连续性。

\textbf{动态安全距离：}冲突判定采用与空域平均速度$\bar{v}$（km/h）正相关的动态阈值：

\begin{equation}
    d_{\text{safety}} = d_{\text{base}} + \bar{v} \times 0.5
\end{equation}

其中$d_{\text{base}} = 50$ m为基础安全距离。当$\exists i \neq j,\ \|(x_i, y_i) - (x_j, y_j)\|_2 < d_{\text{safety}}$时判定为一次冲突。

\subsection{分层奖励函数}

单飞行器即时奖励由四个分量构成：

\begin{equation}
    r_i = r_{\text{safe}} + r_{\text{speed}} + r_{\text{coop}} + r_{\text{eff}}
\end{equation}

\textbf{安全奖励} $r_{\text{safe}}$：三段式分段函数：
\begin{equation}
    r_{\text{safe}} =
    \begin{cases}
        -50 \exp\left(2\frac{d_{\text{safety}} - d_{\min}}{d_{\text{safety}}}\right), & d_{\min} < d_{\text{safety}} \\[8pt]
        -10\frac{1.2d_{\text{safety}} - d_{\min}}{0.2d_{\text{safety}}}, & d_{\text{safety}} \leq d_{\min} < 1.2d_{\text{safety}} \\[8pt]
        5, & d_{\min} \geq 1.2d_{\text{safety}}
    \end{cases}
\end{equation}

\textbf{速度奖励} $r_{\text{speed}}$：理想区间$[65, 90]$ km/h给予正奖励，与速度范围$[60, 140]$ km/h和目标速度$70\pm10$ km/h保持一致：
\begin{equation}
    r_{\text{speed}} =
    \begin{cases}
        -50 \cdot \frac{v_i - 90}{50}, & v_i > 90 \\[8pt]
        -30 \cdot \frac{65 - v_i}{5}, & v_i < 65 \\[8pt]
        20, & 65 \leq v_i \leq 90
    \end{cases}
\end{equation}

\textbf{协同奖励} $r_{\text{coop}}$：$r_{\text{coop}} = -0.2 \cdot |v_i - \bar{v}_{\text{neighbor}}|$，鼓励邻域内速度一致性。

\textbf{效率奖励} $r_{\text{eff}}$：$r_{\text{eff}} = -0.5 \cdot |v_i - v_i^{\text{target}}|$。

全局奖励$R = \frac{1}{N}\sum_i r_i$。当所有$N$架飞行器均进入目标区域（距原点$<10$ m）时，额外给予$+200$的episode完成奖励。

\subsection{动态图注意力网络（DyGAT）}

DyGAT融合三种注意力机制：空间注意力建模同时间步邻域关系，时序注意力捕捉历史轨迹演化趋势，基于物理量的动态边权重区分交互对的风险等级。

\subsubsection{空间注意力}

以LSTM编码特征$h_i^{\text{traj}} \in \mathbb{R}^{64}$为节点特征（$F_{\text{in}}=64$）：

\begin{equation}
    e_{ij} = \text{LeakyReLU}_{0.2}\left(a^{\top}[W h_i^{\text{traj}} \parallel W h_j^{\text{traj}}]\right)
\end{equation}
\begin{equation}
    \alpha_{ij} = \frac{\exp(e_{ij})}{\sum_{k \in \mathcal{N}_i} \exp(e_{ik})},\quad
    h_i^{\text{spatial}} = \sum_{j \in \mathcal{N}_i} \alpha_{ij} \cdot W h_j^{\text{traj}}
\end{equation}

其中$a \in \mathbb{R}^{2F_{\text{out}}}$，$W \in \mathbb{R}^{F_{\text{out}} \times 64}$，$F_{\text{out}}=256$，邻域半径500 m。

\subsubsection{时序注意力}

对最近$T=10$步的历史特征序列计算时序注意力权重$h_i^{\text{temporal}}$。空间与时序特征融合：
\begin{equation}
    h'_i = h_i^{\text{spatial}} + \alpha_{\text{temp}} \cdot h_i^{\text{temporal}}
\end{equation}
其中$\alpha_{\text{temp}} = 0.3$（经验设定，实验中测试了$\{0.1,0.2,0.3,0.5,0.7\}$，0.3在验证集上表现最优）。

\subsubsection{基于物理量的动态边权重}

DyGAT通过基于物理量直接计算的动态边权重编码空间交互的先验知识——距离越近、速度越接近、航向越一致的飞行器对，交互强度越高：

\begin{equation}
    w_{ij} = \exp\left(-\frac{d_{ij}}{1000}\right) \cdot \exp\left(-\frac{|v_i - v_j|}{20}\right) \cdot \exp\left(-\frac{|\theta_i - \theta_j|}{\pi/4}\right)
\end{equation}

三个指数核的函数形式和衰减常数为手工设计，权重值在$(0,1]$之间连续变化。$w_{ij}$通过Hadamard积引导空间注意力：
\begin{equation}
    e_{ij}^{\text{final}} = e_{ij} \odot w_{ij}
\end{equation}

需要指出：三个指数核的衰减常数（1000 m、20 km/h、$\pi/4$ rad）为手工预设而非从数据中学习。这一"物理先验$\times$学习注意力"的双层设计的优势在于：物理先验引导模型将注意力聚焦于高风险交互对，学习注意力在此基础上进行自适应调整；物理先验计算无需可学习参数，不增加梯度回传开销。

\subsection{单智能体网络与GAT-Mixer}

\textbf{单智能体网络}包含四个组件：（1）2层LSTM编码器（64维，Dropout 0.1）；（2）DyGAT层（输出256维$h_i^{\text{coupled}}$）；（3）策略头——2层MLP（256$\to$256$\to$1，Tanh），输出$\mu_i$，实际加速度从$\mathcal{N}(\mu_i, 0.1^2)$采样；（4）Q值头——3层MLP（256$\to$256$\to$128$\to$1，ReLU）。

\textbf{GAT-Mixer混合网络}用于CTDE训练阶段的全局价值估计。将各智能体的4维局部状态$(x_i, y_i, v_i, v_i^{\text{target}})$作为图节点，在全连接图上通过GAT聚合，取节点均值后经MLP输出$Q_{\text{tot}}$。GAT-Mixer的注意力权重不受非负约束，可表达非单调的智能体间价值关系。

\subsection{训练机制}

\subsubsection{课程学习}

训练难度从$N=5$架、$d_{\text{base}}=80$ m开始，在连续三个评估周期满足条件后进阶：$N \leftarrow \min(N+2, 21)$，$d_{\text{base}} \leftarrow \max(d_{\text{base}}-5, 50)$。设有回退机制。最终覆盖$N \in \{5,7,9,11,13,15,17,19,21\}$共九个级别。

\subsubsection{Episode级冲突损失加权}

本文严格保持在PPO的on-policy框架内。PPO的标准裁剪代理目标为：

\begin{equation}
    L^{\text{CLIP}}(\theta) = \hat{\mathbb{E}}_t\left[\min\left(r_t(\theta)\hat{A}_t,\ \text{clip}(r_t(\theta), 1-\epsilon, 1+\epsilon)\hat{A}_t\right)\right]
\end{equation}

其中$r_t(\theta) = \pi_\theta(a_t|o_t) / \pi_{\theta_{\text{old}}}(a_t|o_t)$，$\epsilon = 0.2$。

在此基础上，对包含冲突事件的episode施加更高的损失权重：

\begin{equation}
    L_{\text{total}} = \lambda \cdot (L_{\text{policy}} + L_q + 0.5 \cdot L_{\text{value}})
\end{equation}
\begin{equation}
    \lambda = \begin{cases}
        2.0, & \text{episode内存在冲突事件} \\
        1.0, & \text{无冲突}
    \end{cases}
\end{equation}

需要说明该机制与优先经验回放（PER）的本质区别：PER通过改变采样分布使高TD-error的transition被更频繁地回放，这引入了off-policy数据并需要重要性采样修正；本文的$\lambda$加权仅缩放损失函数的幅度（$\nabla_\theta(\lambda L) = \lambda\nabla_\theta L$），不改变梯度方向，不涉及跨策略版本的数据复用，因此保持在on-policy框架内。该机制的效果等价于在mini-batch中对冲突episode赋予2倍的梯度贡献。

$\lambda=2.0$为在$\{1.0, 1.5, 2.0, 2.5, 3.0\}$五个离散取值中实验确定的最优值（见表\ref{tab:lambda}）。需要指出：五个离散取值中的最优未必是连续空间中的全局最优；在RL训练的高方差条件下，最优值的精确位置存在不确定性。

\begin{table}[H]
    \centering
    \caption{$\lambda$取值的实验结果（21架eVTOL，5次独立测试，均值$\pm$标准差）}
    \label{tab:lambda}
    \begin{tabular}{|c|c|c|c|}
        \hline
        $\lambda$ & TC$\downarrow$ & CFCR(\%)$\uparrow$ & 收敛轮次 \\
        \hline
        1.0 & $24.6 \pm 2.8$ & $53.2 \pm 2.3$ & $\sim800$ \\
        1.5 & $21.8 \pm 2.5$ & $58.7 \pm 2.0$ & $\sim600$ \\
        \textbf{2.0} & $\mathbf{18.2 \pm 2.1}$ & $\mathbf{68.4 \pm 1.8}$ & $\sim400$ \\
        2.5 & $19.5 \pm 2.3$ & $64.1 \pm 1.9$ & $\sim350$ \\
        3.0 & $22.3 \pm 3.1$ & $57.3 \pm 2.5$ & $\sim300$ \\
        \hline
    \end{tabular}
\end{table}

$\lambda > 2.0$时，收敛加快但最终性能下降，推测原因为过度加权导致策略对早期冲突episode过拟合。PTR-PPO\cite{ref27}等工作中也探索了轨迹级优先级在on-policy算法中的应用，采用了更复杂的截断重要性采样机制——本文的$\lambda$加权可视为该方向上一个更为保守和简洁的尝试。

%==========================================================================
% 4. 实验
%==========================================================================
\section{实验与结果分析}

\subsection{实验设置}

实验基于PyTorch实现，硬件环境为NVIDIA GeForce RTX 4060 GPU（8 GB显存）与Intel Core i9-14900K CPU，软件环境为Python 3.11。仿真参数汇总于表\ref{tab:env}。

\begin{table}[H]
    \centering\small
    \caption{仿真环境参数}
    \label{tab:env}
    \begin{tabular}{|c|c|c|c|}
        \hline
        \textbf{参数} & \textbf{取值} & \textbf{参数} & \textbf{取值} \\
        \hline
        空域大小 & $1000\text{m}\times1000\text{m}$ & 最大转弯角 & $\pi/6$ rad \\
        $d_{\text{base}}$ & $50\text{m}$（初始80m） & 时间窗口$T$ & $10$步（5 s） \\
        $\Delta t$ & $0.5$ s & 邻域半径 & $500$ m \\
        最大步长 & $200$步（100 s） & 训练轮次 & $1000$ \\
        目标速度 & $70\pm10$ km/h & $\gamma$ & $0.99$ \\
        速度范围 & $[60, 140]$ km/h & PPO $\epsilon$ & $0.2$ \\
        最大加速度 & $3.0$ m/s$^2$ & 初始学习率$\eta$ & $5\times10^{-4}$ \\
        最大减速度 & $4.0$ m/s$^2$ & 学习率衰减 & $0.99$/轮 \\
        风扰强度 & $0.1$ m/s & 传感器噪声$\sigma$ & $0.5$ m \\
        \hline
    \end{tabular}
\end{table}

\subsubsection{基线方法}

选取八种基线方法，分为四个类别：

\textbf{（1）传统规则调度：}RR（固定轮询）和FCFS（先到先服务）。

\textbf{（2）单/独立智能体RL：}集中式DQN和独立A2C。

\textbf{（3）CTDE-MARL：}MADDPG、COMA、QMIX和MAPPO。

\textbf{（4）工程安全方法：}ORCA\cite{ref4}，通过速度空间几何约束实现分布式碰撞避免。MPC方法在本文场景中存在计算实时性挑战（21架飞行器场景下的单步求解时间通常超过仿真步长500 ms），因此未纳入主实验对比。

\subsubsection{评价指标}

\textbf{表层指标：}累计奖励（CR）、总冲突次数（TC）、平均延迟（AD）。

\textbf{深层指标：}通行完成率（PCR）、无冲突完成率（CFCR）——全程零冲突且所有飞行器完成通行的episode占比（CFCR是PCR的严格子集）、每架次平均冲突次数（MCAE）、最小间隔违反率（MSVR）。

PCR与CFCR的区分是本文的重要方法论贡献：PCR $=99.1\%$与CFCR $=68.4\%$之间约30个百分点的差距揭示了仅报告PCR可能造成的"系统近乎完美"的误导性印象。

\subsection{综合性能对比}

$N=21$架eVTOL测试场景下的综合性能对比如表\ref{tab:comparison}所示（5次独立测试，均值$\pm$标准差）。

\begin{table}[H]
    \centering
    \caption{综合性能对比（21架eVTOL，5次独立测试）}
    \label{tab:comparison}
    \begin{tabular}{|c|c|c|c|c|}
        \hline
        \textbf{方法} & \textbf{CR$\uparrow$} & \textbf{TC$\downarrow$} & \textbf{AD(s)$\downarrow$} & \textbf{PCR(\%)$\uparrow$} \\
        \hline
        RR & $-1256.3\pm89.2$ & $42.6\pm5.3$ & $1.58\pm0.12$ & $89.2\pm3.5$ \\
        FCFS & $-1187.5\pm76.4$ & $38.9\pm4.7$ & $1.42\pm0.10$ & $91.5\pm2.8$ \\
        DQN & $-892.5\pm65.3$ & $28.3\pm3.9$ & $1.21\pm0.09$ & $94.7\pm2.1$ \\
        ORCA & $-198.7\pm35.8$ & $22.8\pm2.9$ & $0.76\pm0.06$ & $96.0\pm1.6$ \\
        独立A2C & $-657.8\pm52.6$ & $19.7\pm2.8$ & $0.93\pm0.07$ & $96.3\pm1.7$ \\
        MADDPG & $-215.6\pm41.2$ & $25.4\pm3.2$ & $0.87\pm0.06$ & $95.8\pm1.9$ \\
        COMA & $187.3\pm35.7$ & $22.1\pm2.9$ & $0.79\pm0.06$ & $97.1\pm1.5$ \\
        QMIX & $426.5\pm28.9$ & $23.1\pm2.7$ & $0.76\pm0.05$ & $97.5\pm1.3$ \\
        MAPPO & $892.4\pm23.5$ & $20.3\pm2.4$ & $0.65\pm0.04$ & $98.2\pm1.1$ \\
        \textbf{本文} & $\mathbf{1286.7\pm18.2}$ & $\mathbf{18.2\pm2.1}$ & $\mathbf{0.57\pm0.03}$ & $\mathbf{99.1\pm0.8}$ \\
        \hline
    \end{tabular}
\end{table}

\textbf{学习方法与规则方法的差异：}本文方法相比RR，TC降低57.2\%（42.6$\to$18.2），AD缩短63.9\%（1.58$\to$0.57 s）。规则方法性能受限于固定调度逻辑无法感知飞行器间相对运动状态。

\textbf{工程安全方法的定位：}ORCA的TC为22.8，显著优于RR和FCFS，验证了分布式几何避障的有效性。但其CFCR仅为约50\%（见表\ref{tab:safety}），低于本文方法，主要原因在于ORCA缺乏对未来多步轨迹的预测能力，高密度场景下易出现避障震荡。

\textbf{与MAPPO的差异：}本文方法相比MAPPO，TC降低10.3\%（20.3$\to$18.2），AD缩短12.3\%。改进幅度在数值上是温和的——MAPPO本身已具备相当的冲突规避能力，本文方法的增量主要来自DyGAT的细粒度时空建模。

\subsubsection{安全指标专项对比}

\begin{table}[H]
    \centering
    \caption{安全指标对比（21架eVTOL，5次独立测试）}
    \label{tab:safety}
    \begin{tabular}{|c|c|c|c|c|}
        \hline
        \textbf{方法} & \textbf{CFCR(\%)$\uparrow$} & \textbf{MCAE$\downarrow$} & \textbf{MSVR(\%)$\downarrow$} & \textbf{PCR(\%)$\uparrow$} \\
        \hline
        RR & $12.3\pm4.1$ & $2.03\pm0.25$ & $8.7\pm1.2$ & $89.2\pm3.5$ \\
        FCFS & $18.7\pm3.6$ & $1.85\pm0.22$ & $7.4\pm1.0$ & $91.5\pm2.8$ \\
        ORCA & $50.2\pm3.0$ & $1.09\pm0.14$ & $3.8\pm0.6$ & $96.0\pm1.6$ \\
        MAPPO & $57.8\pm2.2$ & $0.97\pm0.11$ & $3.0\pm0.4$ & $98.2\pm1.1$ \\
        \textbf{本文} & $\mathbf{68.4\pm1.8}$ & $\mathbf{0.87\pm0.10}$ & $\mathbf{2.4\pm0.3}$ & $\mathbf{99.1\pm0.8}$ \\
        \hline
    \end{tabular}
\end{table}

表\ref{tab:safety}揭示了PCR与CFCR之间约30个百分点的差距——约31\%的episode中，飞行器虽最终到达目标区域但途中至少经历了一次安全距离违反。MAPPO的PCR（98.2\%）与CFCR（57.8\%）之间也存在约40个百分点的差距，表明这一问题是CTDE-MARL方法的共性挑战，并非本文方法特有。

\subsubsection{冲突严重程度分级}

为分析CFCR $=68.4\%$中剩余约31\%存在冲突的episode的严重程度，将安全距离违反按$d_{\min}/d_{\text{safety}}$比值分为三级：

\begin{itemize}
    \item \textbf{轻微}（$0.8 \leq d_{\min}/d_{\text{safety}} < 1.0$）：瞬时小幅偏离，多由噪声或风扰引起。
    \item \textbf{中等}（$0.5 \leq d_{\min}/d_{\text{safety}} < 0.8$）：显著违反，存在真实碰撞风险。
    \item \textbf{严重/NMAC}（$d_{\min}/d_{\text{safety}} < 0.5$）：接近空中碰撞级别，在任何运营标准下均不可接受。
\end{itemize}

分级统计（50次测试episode）表明：约80\%的违反事件属于轻微级别，NMAC级别（$d_{\min} < 42.5$ m）事件占比低于1\%。

\subsection{统计显著性检验}

对本文方法与MAPPO在四项指标上的差异进行Welch双样本$t$检验（双侧，$\alpha=0.05$，$n_1=n_2=5$），结果如表\ref{tab:ttest}所示。

\begin{table}[H]
    \centering
    \caption{本文 vs.\ MAPPO统计检验}
    \label{tab:ttest}
    \begin{tabular}{|c|c|c|c|c|c|}
        \hline
        \textbf{指标} & \textbf{均值差} & \textbf{$t$(df=8)} & \textbf{$p$} & \textbf{Cohen's $d$} & \textbf{结论} \\
        \hline
        TC & $-2.1$ & $1.47$ & $0.180$ & $0.93$ & 不显著 \\
        CFCR & $+10.6\%$ & $8.34$ & $<0.001$ & $5.27$ & 高度显著 \\
        MCAE & $-0.10$ & $1.50$ & $0.172$ & $0.95$ & 不显著 \\
        MSVR & $-0.6\%$ & $2.68$ & $0.029$ & $1.69$ & 显著 \\
        \hline
    \end{tabular}
\end{table}

\textbf{统计效力分析：}基于$n=5$的样本量和Cohen's $d=0.93$的效应量，TC检验的统计效力约为38\%——即在$\alpha=0.05$水平上，即使TC的真实差异存在，当前样本量仅有约38\%的概率检测到。要达到80\%的统计效力，每组需要$n \geq 20$次独立运行。因此，TC的"不显著"应理解为"在$n=5$条件下未检测到显著差异"，而非"确认TC无差异"的证据。相比之下，CFCR的效应量极大（$d=5.27$），即使$n=5$也达到$>99\%$的统计效力，结论高度稳健。MSVR的$d=1.69$（大效应量）在$n=5$下的效力约70\%。

本文方法在CFCR和MSVR上的显著改善——而非TC——构成核心统计贡献。TC检验的统计效力不足是由于样本量限制，而非方法本身缺乏效果（$d=0.93$属于大效应量）。

\subsection{消融实验}

\begin{table}[H]
    \centering
    \caption{消融实验结果（21架eVTOL，$n=5$）}
    \label{tab:ablation}
    \resizebox{0.92\textwidth}{!}{%
    \begin{tabular}{|c|c|c|c|c|}
        \hline
        \textbf{配置} & \textbf{CR$\uparrow$} & \textbf{TC$\downarrow$} & \textbf{AD(s)$\downarrow$} & \textbf{PCR(\%)$\uparrow$} \\
        \hline
        完整模型 & $\mathbf{1286.7\pm18.2}$ & $\mathbf{18.2\pm2.1}$ & $\mathbf{0.57\pm0.03}$ & $\mathbf{99.1\pm0.8}$ \\
        $-$ DyGAT（标准GAT） & $215.3\pm32.6$ & $49.8\pm4.5$ & $0.79\pm0.05$ & $94.3\pm1.9$ \\
        $-$ 课程学习 & $587.2\pm27.4$ & $28.4\pm3.1$ & $0.68\pm0.04$ & $97.2\pm1.4$ \\
        $-$ 冲突加权（$\lambda\equiv1.0$） & $763.5\pm24.1$ & $24.6\pm2.8$ & $0.65\pm0.04$ & $97.8\pm1.2$ \\
        $-$ 课程学习 $-$ 冲突加权 & $128.6\pm38.5$ & $35.7\pm3.9$ & $0.82\pm0.06$ & $95.6\pm1.7$ \\
        \hline
    \end{tabular}%
    }
\end{table}

DyGAT是性能贡献最大的单一模块，移除后TC增加173.6\%。DyGAT移除的整体替换方式尚未分解空间注意力、时序注意力与动态边权重的独立贡献——更细粒度的组件级消融是进一步验证各子模块独立效果的必要工作。

课程学习与冲突加权的作用路径基本独立：单独移除分别导致TC增加56.0\%和35.2\%，两者同时移除增加96.2\%。

\subsection{密度鲁棒性与敏感性分析}

\textbf{密度鲁棒性：}在$N \in \{5, 10, 15, 21\}$四种密度下，本文方法的TC增长斜率（$\approx0.75$次/架）低于MAPPO（$\approx1.15$次/架）和RR（$\approx2.4$次/架）。DyGAT的动态边权重在高密度下优势更为显著。

\textbf{参数敏感性：}$d_{\text{base}}$从40 m增至70 m时，TC下降34.7\%但AD上升21.2\%，呈现经典的安全-效率权衡。风扰强度从0增至0.3 m/s时TC上升约28\%，传感器噪声$\sigma$从0增至1.5 m时TC上升约15\%。

\subsection{推理效率}

本文方法在$N=21$时的单步推理时间为4.5 ms（RTX 4060），实时性裕度约为$111\times$（相对500 ms仿真步长）。计算时间增长斜率约0.21 ms/架，在所有CTDE-MARL方法中最低。

%==========================================================================
% 5. 讨论
%==========================================================================
\section{讨论}

\subsection{与MS-GRL的比较}

Li等\cite{ref10}提出的MS-GRL与本文方法在三个层面存在本质差异：

\textbf{算法选择：}MS-GRL采用Double DQN（离散速度档位），训练稳定但控制粒度粗糙；本文采用PPO（连续加速度），适合eVTOL的精细化速度调控需求。

\textbf{图结构设计：}MS-GRL的多尺度GAT基于固定距离阈值。本文DyGAT通过距离、速度差和航向差三个物理量计算连续动态边权重$w_{ij}\in(0,1]$，可区分"距离相同但相对运动状态不同"的飞行器对的差异化风险——例如两对均为200 m间距的飞行器，同向等速（$w_{ij}\approx1.0$）与相向高速靠近（$w_{ij}\approx0.015$）的注意力权重相差约67倍。

\textbf{场景针对性：}MS-GRL面向100架通用UAV场景，本文聚焦21架eVTOL交叉路口的特定运营场景。

两者在图增强MARL技术路线上互补。需要指出的是，MS-GRL论文在IEEE IoT-J发表时的实际作者为Li Y, Li J, Wang J, Zhang X, Ding H, Du W，卷期页码为12(21):44290--44303\cite{ref10}。

\subsection{GAT-Mixer的设计说明}

本文的GAT-Mixer通过GAT注意力聚合和均值池化实现全局Q值估计。需要说明：ICLR 2023发表的GraphMixer\cite{ref28}的核心设计为基于MLP-Mixer的link-encoder/node-encoder/mean-pooling架构，本文的GAT-Mixer仅借用了"Mixer"的命名（指混合多个智能体的Q值），在技术路线上与GraphMixer无关——GAT-Mixer更接近一个"基于GAT编码器的QMIX变体"。这一命名上的关联性在严格意义上不够精确，在此予以澄清。

\subsection{局限性分析}

本文方法存在以下主要局限性：

\textbf{（1）场景维度。}当前仿真为二维单交叉路口场景。真实低空交通涉及三维高度层、多路口网络拓扑、起降场容量约束等维度。二维简化是当前工作最根本的场景限制。

\textbf{（2）统计样本量。}$n=5$的独立测试次数的统计效力有限。对TC等效应量中等（$d\approx0.93$）的指标，当前样本量的检验效力仅约38\%。扩大样本量至$n\geq20$以提供更可靠的统计推断是必要的后续工作。

\textbf{（3）风扰与噪声模型。}当前采用均匀随机风扰和高斯白噪声，均为时间独立的简化模型。真实大气边界层风场具有显著的空间相关性和时间连续性（Dryden/Kaimal谱模型），城市环境中还存在建筑尾流等复杂效应。

\textbf{（4）通信假设。}CTDE训练假设全局状态可无延迟获取。实际UTM部署中通信链路存在延迟（50--200 ms）和丢包约束。

\textbf{（5）eVTOL动力学参数。}第3.2.2节使用的加速度取值（3.0--4.0 m/s$^2$）参考了学术文献中对eVTOL的典型建模假设，但当前主流制造商并未公开确切的最大水平加速度规格。在更精确的制造商数据可用时，应对这些参数进行更新和验证。

\textbf{（6）基线覆盖。}MPC和CBF两种具有形式化安全保证的方法尚未纳入实验对比。

%==========================================================================
% 6. 结论
%==========================================================================
\section{结论}

本文针对城市低空十字交叉路口的eVTOL协同管控问题，提出了一种基于DyGAT的耦合MARL模型。在21架eVTOL的二维交叉路口仿真场景中，本文方法在CFCR（68.4\% vs.\ MAPPO的57.8\%，$p<0.001$）和MSVR（2.4\% vs.\ 3.0\%，$p=0.029$）上取得了统计显著的改善。TC的改善（18.2 vs.\ 20.3，降幅10.3\%）具有大效应量（$d=0.93$），但在$n=5$的样本量下未达统计显著（$p=0.180$）——这需要更大规模的重复实验来验证。

本文的方法论贡献在于：（1）物理量驱动的动态边权重为图结构MARL中的时空耦合建模提供了一种简洁有效的方法；（2）区分表层指标与深层指标的评价框架揭示了仅报告PCR可能造成的误导性印象；（3）系统性的统计效力分析诚实呈现了当前实证证据的强度与局限。

后续工作方向包括：（1）将仿真场景扩展至三维多交叉路口网络；（2）在更大计算资源下进行$n\geq20$的独立重复实验；（3）集成标准风谱模型（Dryden/Kaimal）替代当前简化风扰；（4）引入通信延迟与丢包约束评估分布式执行鲁棒性；（5）与MPC、CBF等安全关键方法进行系统基准对比。

%==========================================================================
% 参考文献
%==========================================================================
\begin{thebibliography}{99}

\bibitem{ref1} Federal Aviation Administration. Urban air mobility (UAM): concept of operations v2.0[R]. Washington, DC: FAA, 2023.

\bibitem{ref2} MENON P K, SWERIDUK G D, BILIMORIA K D. New approach to modeling, analysis, and control of air traffic flow[J]. Journal of Guidance, Control, and Dynamics, 2004, 27(5): 737--744.

\bibitem{ref3} RICHARDS A, HOW J P. Aircraft trajectory planning with collision avoidance using mixed integer linear programming[C]//Proceedings of the 2002 American Control Conference. Piscataway: IEEE, 2002: 1936--1941.

\bibitem{ref4} VAN DEN BERG J, GUY S J, LIN M, et al. Reciprocal n-body collision avoidance[C]//Robotics Research: The 14th International Symposium ISRR (Springer Tracts in Advanced Robotics, Vol.\ 70). Berlin: Springer, 2011: 3--19.

\bibitem{ref5} AMES A D, COOGAN S, EGERSTEDT M, et al. Control barrier functions: theory and applications[C]//2019 18th European Control Conference (ECC). Piscataway: IEEE, 2019: 3420--3431.

\bibitem{ref6} AGOGINO A, TUMER K. Regulating air traffic flow with coupled agents[C]//Proceedings of the 7th International Conference on Autonomous Agents and Multiagent Systems (AAMAS 2008). Richland: IFAAMAS, 2008: 535--542.

\bibitem{ref7} AGOGINO A K, TUMER K. A multiagent approach to managing air traffic flow[J]. Autonomous Agents and Multi-Agent Systems, 2012, 24(1): 1--25.

\bibitem{ref8} BRITTAIN M, WEI P. Autonomous separation assurance in a high-density en route sector: a deep multi-agent reinforcement learning approach[C]//2019 IEEE Intelligent Transportation Systems Conference (ITSC). Piscataway: IEEE, 2019: 3256--3262.

\bibitem{ref9} GHOSH S, LAGUNA S, LIM S H, et al. A deep ensemble method for multi-agent reinforcement learning: a case study on air traffic control[C]//Proceedings of the 31st International Conference on Automated Planning and Scheduling (ICAPS 2021). Palo Alto: AAAI Press, 2021: 468--476.

\bibitem{ref10} LI Y, LI J, WANG J, ZHANG X, DING H, DU W. Multiscale-graph-enhanced reinforcement learning for conflict resolution in dense UAV networks[J]. IEEE Internet of Things Journal, 2025, 12(21): 44290--44303.

\bibitem{ref11} ASLANI M, MESGARI M S, WIERING M. Adaptive traffic signal control with actor-critic methods in a real-world traffic network with different traffic disruption events[J]. Transportation Research Part C: Emerging Technologies, 2017, 85: 732--752.

\bibitem{ref12} 翟子洋, 郝茹茹, 董世浩. 大规模智慧交通信号控制中的强化学习和深度强化学习方法综述[J]. 计算机应用研究, 2024, 41(6): 1618--1627.

\bibitem{ref13} SPATHARIS C, BASTAS A, KRAVARIS T, et al. Hierarchical multiagent reinforcement learning schemes for air traffic management[J]. Neural Computing and Applications, 2021, 33(14): 8649--8671.

\bibitem{ref14} TUMER K, AGOGINO A. Distributed agent-based air traffic flow management[C]//Proceedings of the 6th International Joint Conference on Autonomous Agents and Multiagent Systems (AAMAS 2007). Honolulu: ACM, 2007: 330--337.

\bibitem{ref15} 王迎. 城市交通信号控制深度强化学习算法研究[D]. 济南: 山东交通学院, 2024.

\bibitem{ref16} 赵晓华, 石建军, 李振龙, 赵国勇. 基于Q-learning和BP神经元网络的交叉口信号灯控制[J]. 公路交通科技, 2007, 24(7): 99--102.

\bibitem{ref17} 赖建辉. 基于深度强化学习的交通控制优化方法研究及实现[D]. 杭州: 浙江工业大学, 2020.

\bibitem{ref18} 张新武. 基于深度强化学习算法的交通信号控制优化研究[D]. 太原: 太原科技大学, 2024.

\bibitem{ref19} 承向军, 常歆识, 杨肇夏. 基于Q-学习的交通信号控制方法[J]. 系统工程理论与实践, 2006(8): 136--140.

\bibitem{ref20} 卢守峰, 张术, 刘喜敏. 平均排队长度差最小的单交叉口在线Q学习模型[J]. 公路交通科技, 2014, 31(11): 116--122.

\bibitem{ref21} 江波, 李彦冬, 刘威陇, 等. 应用深度Q学习的机场道口协同智能调速方法[J]. 航空计算技术, 2022, 52(1): 26--30.

\bibitem{ref22} 张春阳, 席雅琪. 人工智能在城市交通控制中的应用综述[J]. 信息系统工程, 2024(07): 33--36.

\bibitem{ref23} GHAVAMZADEH M, MAHADEVAN S, MAKAR R. Hierarchical multi-agent reinforcement learning[J]. Autonomous Agents and Multi-Agent Systems, 2006, 13(2): 197--229.

\bibitem{ref24} TUMER K, AGOGINO A. Improving air traffic management with a learning multiagent system[J]. IEEE Intelligent Systems, 2009, 24(1): 18--21.

\bibitem{ref25} BACCHINI A, CESTINO E. Electric VTOL configurations comparison[J]. Aerospace, 2019, 6(3): 26.

\bibitem{ref26} U.S. Department of Defense. Flying qualities of piloted airplanes: MIL-F-8785C[S]. Washington, DC: U.S. Department of Defense, 1980.

\bibitem{ref27} LIANG X, MA Y, FENG Y, et al. PTR-PPO: proximal policy optimization with prioritized trajectory replay[EB/OL]. arXiv:2112.03798, 2021.

\bibitem{ref28} CONG W, ZHANG S, CHEN J, et al. Do we really need complicated model architectures for temporal networks?[C]//Proceedings of the 11th International Conference on Learning Representations (ICLR 2023). Kigali: ICLR, 2023.

\bibitem{ref29} XIE Y B, GARDI A, SABATINI R. Reinforcement learning-based flow management techniques for urban air mobility and dense low-altitude air traffic operations[C]//2021 IEEE/AIAA 40th Digital Avionics Systems Conference (DASC). Piscataway: IEEE, 2021: 1--10.

\end{thebibliography}

\end{document}
